% GENERATED FILE. Edit canonical lessons, then run npm run generate:books.
% canonical-source-hash: fnv1a64:346a5f7d580b5419
% canonical-lessons: UR-C469-dakar, UR-C469-hichki, UR-C469-khankarna, UR-C469-balgham, UR-C469-nazla

\chapter{Hiccups and Colds}
\label{ch:ur-hiccups-and-colds}
\hlchaptermodality{469}

\begin{chapteropening}
\textbf{By the end of this chapter:} \emph{I can talk about hiccups, coughs and colds.}

Complete the last lesson of chapter 469: I can talk about hiccups, coughs and colds.
\end{chapteropening}

\section[\texorpdfstring{\ur{ڈکار} (ḍakār)}{ڈکار (ḍakār)}]{\ur{ڈکار} (ḍakār) — a burp}

\label{lesson:UR-C469-dakar}



\begin{quote}
Before the new one: say the Urdu for to scratch an itch, then the Urdu for a rash.
\end{quote}

\subsection*{You'll want to know: \ur{ڈکار}}
\textbf{\ur{ڈکار}} — \emph{ḍakār} — \textquotedblleft{}a burp\textquotedblright{}.

\begin{quote}
\ur{کھانے} \ur{کے} \ur{بعد} \ur{اسے} \ur{ڈکار} \ur{آئی۔}
\end{quote}

\begin{quote}
\emph{khāne ke bād use ḍakār āī.}
\end{quote}

\begin{quote}
He burped after the meal.
\end{quote}

Feminine.

\begin{cousinweb}
\textbf{\ur{ڈکار}} (\emph{ḍakār}) is inherited from \emph{ḍakkāra}, \textquotedblleft{}belch\textquotedblright{}, a form CDIAL reconstructs and treats as an imitation of the sound (CDIAL 5521).
\end{cousinweb}

\begin{grammarlens}[title={What you've built}]
One more word for describing what you lived and felt.
\end{grammarlens}

\subsection*{Your turn}
\begin{itemize}
\raggedright
  \item \emph{Say it:} \emph{ḍakār}
  \item \emph{Say it:} \emph{ḍakār}, once more
  \item \emph{Say it:} \emph{khāne ke bād use ḍakār āī.} — He burped after the meal.
  \item \emph{Recall:} say the Urdu for to scratch an itch, then the Urdu for a rash, then say \emph{ḍakār} again
\end{itemize}

\begin{tcolorbox}[breakable,colback=teal!4,colframe=teal!35!black,title={Before you move on}]
What does \ur{ڈکار} mean? (\textquotedblleft{}A burp\textquotedblright{}.) Say it once more.
\end{tcolorbox}
\section[\texorpdfstring{\ur{ہچکی} (hichkī)}{ہچکی (hichkī)}]{\ur{ہچکی} (hichkī) — a hiccup}

\label{lesson:UR-C469-hichki}



\begin{quote}
Before the new one: say the Urdu for a rash, then the Urdu for a burp.
\end{quote}

\subsection*{You'll want to know: \ur{ہچکی}}
\textbf{\ur{ہچکی}} — \emph{hichkī} — \textquotedblleft{}a hiccup\textquotedblright{}.

\begin{quote}
\ur{پانی} \ur{پینے} \ur{سے} \ur{ہچکی} \ur{رک} \ur{گئی۔}
\end{quote}

\begin{quote}
\emph{pānī pīne se hichkī ruk gaī.}
\end{quote}

\begin{quote}
Drinking some water stopped the hiccups.
\end{quote}

Feminine.

\begin{cousinweb}
is inherited: CDIAL lists the Hindi \emph{hīk}, \textquotedblleft{}hiccup\textquotedblright{}, under Sanskrit \emph{hikkā}, \textquotedblleft{}hiccup\textquotedblright{} (CDIAL 14075); \textbf{\ur{ہچکی}} is probably a longer form of the same imitative word.
\end{cousinweb}

\begin{grammarlens}[title={What you've built}]
One more word for describing what you lived and felt.
\end{grammarlens}

\subsection*{Your turn}
\begin{itemize}
\raggedright
  \item \emph{Say it:} \emph{hichkī}
  \item \emph{Say it:} \emph{hichkī}, once more
  \item \emph{Say it:} \emph{pānī pīne se hichkī ruk gaī.} — Drinking some water stopped the hiccups.
  \item \emph{Recall:} say the Urdu for a rash, then the Urdu for a burp, then say \emph{hichkī} again
\end{itemize}

\begin{tcolorbox}[breakable,colback=teal!4,colframe=teal!35!black,title={Before you move on}]
What does \ur{ہچکی} mean? (\textquotedblleft{}A hiccup\textquotedblright{}.) Say it once more.
\end{tcolorbox}
\section[\texorpdfstring{\ur{کھنکارنا} (khaṅkārnā)}{کھنکارنا (khaṅkārnā)}]{\ur{کھنکارنا} (khaṅkārnā) — to clear one's throat}

\label{lesson:UR-C469-khankarna}



\begin{quote}
Before the new one: say the Urdu for a burp, then the Urdu for a hiccup.
\end{quote}

\subsection*{You'll want to know: \ur{کھنکارنا}}
\textbf{\ur{کھنکارنا}} — \emph{khaṅkārnā} — \textquotedblleft{}to clear one's throat\textquotedblright{}.

\begin{quote}
\ur{بولنے} \ur{سے} \ur{پہلے} \ur{دادا} \ur{نے} \ur{کھنکارا۔}
\end{quote}

\begin{quote}
\emph{bolne se pehle dādā ne khaṅkārā.}
\end{quote}

\begin{quote}
My grandfather cleared his throat before speaking.
\end{quote}

Past, with \ur{نے} (he or she …): \ur{اس} \ur{نے} \ur{کھنکارا}, \emph{us ne khaṅkārā}.

\begin{cousinweb}
\textbf{\ur{کھنکارنا}} (\emph{khaṅkārnā}) is inherited: CDIAL lists the Hindi cousin \emph{khãkārnā} under \emph{khaṅkh}, \textquotedblleft{}cough, hawk\textquotedblright{}, a form it reconstructs and treats as an imitation of the sound (CDIAL 3763).
\end{cousinweb}

\begin{grammarlens}[title={What you've built}]
One more word for describing what you lived and felt.
\end{grammarlens}

\subsection*{Your turn}
\begin{itemize}
\raggedright
  \item \emph{Say it:} \emph{khaṅkārnā}
  \item \emph{Say it:} \emph{khaṅkārnā}, once more
  \item \emph{Say it:} \emph{bolne se pehle dādā ne khaṅkārā.} — My grandfather cleared his throat before speaking.
  \item \emph{Recall:} say the Urdu for a burp, then the Urdu for a hiccup, then say \emph{khaṅkārnā} again
\end{itemize}

\begin{tcolorbox}[breakable,colback=teal!4,colframe=teal!35!black,title={Before you move on}]
What does \ur{کھنکارنا} mean? (\textquotedblleft{}To clear one's throat\textquotedblright{}.) Say it once more.
\end{tcolorbox}
\section[\texorpdfstring{\ur{بلغم} (balgham)}{بلغم (balgham)}]{\ur{بلغم} (balgham) — phlegm}

\label{lesson:UR-C469-balgham}



\begin{quote}
Before the new one: say the Urdu for a hiccup, then the Urdu for to clear one's throat.
\end{quote}

\subsection*{You'll want to know: \ur{بلغم}}
\textbf{\ur{بلغم}} — \emph{balgham} — \textquotedblleft{}phlegm\textquotedblright{}.

\begin{quote}
\ur{کھانسی} \ur{کے} \ur{ساتھ} \ur{بلغم} \ur{آتا} \ur{ہے۔}
\end{quote}

\begin{quote}
\emph{khāṅsī ke sāth balgham ātā hai.}
\end{quote}

\begin{quote}
Phlegm comes up with the cough.
\end{quote}

Masculine.

\begin{cousinweb}
\textbf{\ur{بلغم}} (\emph{balgham}) is usually traced, through Persian, to Arabic \emph{balgham}, \textquotedblleft{}phlegm\textquotedblright{}, which Arabic probably took from Greek \emph{phlegma}.
\end{cousinweb}

\begin{grammarlens}[title={What you've built}]
One more word for describing what you lived and felt.
\end{grammarlens}

\subsection*{Your turn}
\begin{itemize}
\raggedright
  \item \emph{Say it:} \emph{balgham}
  \item \emph{Say it:} \emph{balgham}, once more
  \item \emph{Say it:} \emph{khāṅsī ke sāth balgham ātā hai.} — Phlegm comes up with the cough.
  \item \emph{Recall:} say the Urdu for a hiccup, then the Urdu for to clear one's throat, then say \emph{balgham} again
\end{itemize}

\begin{tcolorbox}[breakable,colback=teal!4,colframe=teal!35!black,title={Before you move on}]
What does \ur{بلغم} mean? (\textquotedblleft{}Phlegm\textquotedblright{}.) Say it once more.
\end{tcolorbox}
\section[\texorpdfstring{\ur{نزلہ} (nazlā)}{نزلہ (nazlā)}]{\ur{نزلہ} (nazlā) — catarrh; a running cold}

\label{lesson:UR-C469-nazla}



\begin{quote}
Before the new one: say the Urdu for to clear one's throat, then the Urdu for phlegm.
\end{quote}

\subsection*{You'll want to know: \ur{نزلہ}}
\textbf{\ur{نزلہ}} — \emph{nazlā} — \textquotedblleft{}catarrh; a running cold\textquotedblright{}.

\begin{quote}
\ur{سردی} \ur{میں} \ur{اسے} \ur{نزلہ} \ur{ہو} \ur{گیا۔}
\end{quote}

\begin{quote}
\emph{sardī meṅ use nazlā ho gayā.}
\end{quote}

\begin{quote}
In the cold weather she caught a streaming cold.
\end{quote}

Masculine.

\begin{cousinweb}
\textbf{\ur{نزلہ}} (\emph{nazlā}) is usually traced, through Persian, to Arabic \emph{nazla}, \textquotedblleft{}a flowing down\textquotedblright{}, and so a running cold.
\end{cousinweb}

\begin{grammarlens}[title={What you've built}]
One more word for describing what you lived and felt.
\end{grammarlens}

\subsection*{Your turn}
\begin{itemize}
\raggedright
  \item \emph{Say it:} \emph{nazlā}
  \item \emph{Say it:} \emph{nazlā}, once more
  \item \emph{Say it:} \emph{sardī meṅ use nazlā ho gayā.} — In the cold weather she caught a streaming cold.
  \item \emph{Recall:} say the Urdu for to clear one's throat, then the Urdu for phlegm, then say \emph{nazlā} again
\end{itemize}

\begin{tcolorbox}[breakable,colback=teal!4,colframe=teal!35!black,title={Before you move on}]
What does \ur{نزلہ} mean? (\textquotedblleft{}Catarrh; a running cold\textquotedblright{}.) Say it once more.
\end{tcolorbox}
